\documentclass[a4paper]{article}
\usepackage[utf8]{inputenc}
\usepackage[T1]{fontenc}
\usepackage[spanish]{babel}
\usepackage{geometry}
\usepackage{tikz}
\usepackage{xcolor}
\usepackage{fontspec}
\usepackage{graphicx}

\geometry{left=5mm, right=5mm, top=5mm, bottom=5mm}
\pagestyle{empty}

% === COLORES ===
% Campo semantico (oscuros)
\definecolor{violeta}{HTML}{6B3FA0}
\definecolor{naranja}{HTML}{D4730E}
% Genero (claros)
\definecolor{azulceleste}{HTML}{7EC8E3}
\definecolor{rojocoral}{HTML}{F28B82}
% Texto
\definecolor{azulmarino}{HTML}{1B3A5C}
\definecolor{grisFondo}{HTML}{F2F2F2}
% Fondo tarjeta vocabulario
\definecolor{rosapastel}{HTML}{F2D7D9}
% Fondo imagen
\definecolor{purpuraImg}{HTML}{7B5EA7}
% Caja blanca semitransparente
\definecolor{cajaBlanca}{HTML}{FFFFFF}

% === TIPOGRAFIA ===
\setmainfont{Montserrat}[
  BoldFont = {Montserrat Bold},
  ItalicFont = {Montserrat Italic},
]

% === DIMENSIONES TARJETA ===
\newlength{\cardW}\setlength{\cardW}{63mm}
\newlength{\cardH}\setlength{\cardH}{88mm}
\newlength{\margin}\setlength{\margin}{3mm}

% === COMANDO TARJETA FRONTAL ===
% #1 = palabra
% #2 = color campo semantico (violeta/naranja)
% #3 = color genero (azulceleste/rojocoral)
% #4 = frecuencia (1-3)
% #5 = frase de ejemplo
\newcommand{\tarjetaFrontal}[5]{%
\begin{tikzpicture}
  % Fondo tarjeta
  \fill[#3, rounded corners=2mm] (0,0) rectangle (\cardW,\cardH);
  % Franja superior (campo semantico)
  \fill[#2, rounded corners=2mm] (0,\cardH-10mm) rectangle (\cardW,\cardH);
  \fill[#2] (0,\cardH-12mm) rectangle (\cardW,\cardH-10mm);
  % Palabra
  \node[anchor=center, text=white, font=\bfseries\Large] at (0.5\cardW, \cardH-6mm) {\MakeUppercase{#1}};
  % Frecuencia (barras)
  \ifnum#4>0
    \fill[white] ({\cardW-10mm},{\cardH-4mm}) rectangle ({\cardW-8mm},{\cardH-8mm});
  \fi
  \ifnum#4>1
    \fill[white] ({\cardW-7mm},{\cardH-3mm}) rectangle ({\cardW-5mm},{\cardH-8mm});
  \fi
  \ifnum#4>2
    \fill[white] ({\cardW-4mm},{\cardH-2mm}) rectangle ({\cardW-2mm},{\cardH-8mm});
  \fi
  % Zona imagen (placeholder)
  \fill[purpuraImg, rounded corners=1mm] (\margin+2mm,\cardH-68mm) rectangle (\cardW-\margin-2mm,\cardH-12mm);
  \node[anchor=center, text=white, font=\small\itshape, text width=40mm, align=center] at (0.5\cardW, \cardH-40mm) {[imagen\\#1]};
  % Caja blanca semitransparente para frase
  \fill[white, opacity=0.85, rounded corners=2mm] (\margin, \margin) rectangle (\cardW-\margin, \cardH-70mm);
  % Frase de ejemplo
  \node[anchor=center, text=azulmarino, font=\small, text width=50mm, align=center] at (0.5\cardW, \cardH-75mm) {#5};
\end{tikzpicture}%
}

% === COMANDO TARJETA POSTERIOR ===
% #1 = palabra
% #2 = color campo semantico
% #3 = color genero
% #4 = combo1
% #5 = combo2
% #6 = combo3
% #7 = combo4
% #8 = traducciones (IT, FR, PT, EN, CS, PL)
\newcommand{\tarjetaPosterior}[8]{%
\begin{tikzpicture}
  % Fondo tarjeta
  \fill[#3, rounded corners=2mm] (0,0) rectangle (\cardW,\cardH);
  % Franja superior
  \fill[#2, rounded corners=2mm] (0,\cardH-10mm) rectangle (\cardW,\cardH);
  \fill[#2] (0,\cardH-12mm) rectangle (\cardW,\cardH-10mm);
  % Palabra
  \node[anchor=center, text=white, font=\bfseries\Large] at (0.5\cardW, \cardH-6mm) {\MakeUppercase{#1}};
  % Caja blanca para combos
  \fill[white, opacity=0.85, rounded corners=2mm] (\margin, \cardH-70mm) rectangle (\cardW-\margin, \cardH-13mm);
  % Titulo Combos
  \node[anchor=north west, text=azulmarino, font=\bfseries\normalsize] at (\margin+2mm, \cardH-15mm) {Combos};
  % Combos
  \node[anchor=north west, text=azulmarino, font=\small, text width=50mm] at (\margin+2mm, \cardH-22mm) {%
    #4\\[1mm]
    #5\\[1mm]
    #6\\[1mm]
    #7%
  };
  % Caja blanca para traducciones
  \fill[white, opacity=0.85, rounded corners=2mm] (\margin, \margin) rectangle (\cardW-\margin, \cardH-72mm);
  % Traducciones
  \node[anchor=north west, text=azulmarino, font=\scriptsize, text width=50mm] at (\margin+2mm, \cardH-74mm) {#8};
\end{tikzpicture}%
}

\begin{document}

%% ============================================================
%% CARAS FRONTALES — Pagina 1 (8 tarjetas)
%% ============================================================

\noindent
\tarjetaFrontal{padre}{violeta}{azulceleste}{3}{El padre de Javier es camionero.}%
\hspace{1mm}%
\tarjetaFrontal{madre}{violeta}{rojocoral}{3}{La madre de Javier se llama Catalina.}%

\vspace{1mm}

\noindent
\tarjetaFrontal{hermano}{violeta}{azulceleste}{3}{Carlos es el hermano mayor de Lucia.}%
\hspace{1mm}%
\tarjetaFrontal{hermana}{violeta}{rojocoral}{3}{Javier tiene una hermana, Alejandra.}%

\vspace{1mm}

\noindent
\tarjetaFrontal{hijo}{violeta}{azulceleste}{2}{David es el hijo de Alicia y Luis.}%
\hspace{1mm}%
\tarjetaFrontal{hija}{violeta}{rojocoral}{2}{Maria es la hija de Carlos y Carmen.}%

\vspace{1mm}

\noindent
\tarjetaFrontal{abuelo}{violeta}{azulceleste}{2}{El abuelo de David se llama Carlos.}%
\hspace{1mm}%
\tarjetaFrontal{abuela}{violeta}{rojocoral}{2}{La abuela de David se llama Carmen.}%

\newpage

%% ============================================================
%% CARAS FRONTALES — Pagina 2 (8 tarjetas)
%% ============================================================

\noindent
\tarjetaFrontal{tio}{violeta}{azulceleste}{2}{Nacho es el tio de David.}%
\hspace{1mm}%
\tarjetaFrontal{tia}{violeta}{rojocoral}{2}{Maria es la tia de David.}%

\vspace{1mm}

\noindent
\tarjetaFrontal{primo}{violeta}{azulceleste}{2}{Alvaro es el primo de David.}%
\hspace{1mm}%
\tarjetaFrontal{prima}{violeta}{rojocoral}{2}{Paloma es la prima de David.}%

\vspace{1mm}

\noindent
\tarjetaFrontal{marido}{violeta}{azulceleste}{2}{Luis es el marido de Alicia.}%
\hspace{1mm}%
\tarjetaFrontal{mujer}{violeta}{rojocoral}{3}{Alicia es la mujer de Luis.}%

\vspace{1mm}

\noindent
\tarjetaFrontal{nieto}{violeta}{azulceleste}{1}{David es el nieto de Carlos y Carmen.}%
\hspace{1mm}%
\tarjetaFrontal{nieta}{violeta}{rojocoral}{1}{Pilar es la nieta de Carlos y Carmen.}%

\newpage

%% ============================================================
%% CARAS FRONTALES — Pagina 3 (5 tarjetas restantes)
%% ============================================================

\noindent
\tarjetaFrontal{sobrino}{violeta}{azulceleste}{1}{David es el sobrino de Roberto.}%
\hspace{1mm}%
\tarjetaFrontal{sobrina}{violeta}{rojocoral}{1}{Paloma es la sobrina de Alicia.}%

\vspace{1mm}

\noindent
\tarjetaFrontal{camionero}{naranja}{azulceleste}{1}{El padre de Javier es camionero.}%
\hspace{1mm}%
\tarjetaFrontal{enfermera}{naranja}{rojocoral}{1}{Catalina es enfermera.}%

\vspace{1mm}

\noindent
\tarjetaFrontal{agricultor}{naranja}{azulceleste}{1}{El padre de Lucia es agricultor.}%

\newpage

%% ============================================================
%% CARAS POSTERIORES — Pagina 4 (8 tarjetas)
%% ============================================================

\noindent
\tarjetaPosterior{padre}{violeta}{azulceleste}%
  {el/la + padre + de + nombre}%
  {mi/tu/su + padre}%
  {es + el padre de + nombre}%
  {padre + madre = padres}%
  {IT: padre | FR: pere | PT: pai | EN: father | CS: otec | PL: ojciec}%
\hspace{1mm}%
\tarjetaPosterior{madre}{violeta}{rojocoral}%
  {el/la + madre + de + nombre}%
  {mi/tu/su + madre}%
  {es + la madre de + nombre}%
  {padre + madre = padres}%
  {IT: madre | FR: mere | PT: mae | EN: mother | CS: matka | PL: matka}%

\vspace{1mm}

\noindent
\tarjetaPosterior{hermano}{violeta}{azulceleste}%
  {el/la + hermano + de + nombre}%
  {mi/tu/su + hermano}%
  {es + el hermano de + nombre}%
  {hermano + hermana = hermanos}%
  {IT: fratello | FR: frere | PT: irmao | EN: brother | CS: bratr | PL: brat}%
\hspace{1mm}%
\tarjetaPosterior{hermana}{violeta}{rojocoral}%
  {el/la + hermana + de + nombre}%
  {mi/tu/su + hermana}%
  {es + la hermana de + nombre}%
  {hermano + hermana = hermanos}%
  {IT: sorella | FR: soeur | PT: irma | EN: sister | CS: sestra | PL: siostra}%

\vspace{1mm}

\noindent
\tarjetaPosterior{hijo}{violeta}{azulceleste}%
  {el/la + hijo + de + nombre}%
  {mi/tu/su + hijo}%
  {es + el hijo de + nombre}%
  {hijo + hija = hijos}%
  {IT: figlio | FR: fils | PT: filho | EN: son | CS: syn | PL: syn}%
\hspace{1mm}%
\tarjetaPosterior{hija}{violeta}{rojocoral}%
  {el/la + hija + de + nombre}%
  {mi/tu/su + hija}%
  {es + la hija de + nombre}%
  {hijo + hija = hijos}%
  {IT: figlia | FR: fille | PT: filha | EN: daughter | CS: dcera | PL: corka}%

\vspace{1mm}

\noindent
\tarjetaPosterior{abuelo}{violeta}{azulceleste}%
  {el/la + abuelo + de + nombre}%
  {mi/tu/su + abuelo}%
  {es + el abuelo de + nombre}%
  {abuelo + abuela = abuelos}%
  {IT: nonno | FR: grand-pere | PT: avo | EN: grandfather | CS: dedecek | PL: dziadek}%
\hspace{1mm}%
\tarjetaPosterior{abuela}{violeta}{rojocoral}%
  {el/la + abuela + de + nombre}%
  {mi/tu/su + abuela}%
  {es + la abuela de + nombre}%
  {abuelo + abuela = abuelos}%
  {IT: nonna | FR: grand-mere | PT: avo | EN: grandmother | CS: babicka | PL: babcia}%

\newpage

%% ============================================================
%% CARAS POSTERIORES — Pagina 5 (8 tarjetas)
%% ============================================================

\noindent
\tarjetaPosterior{tio}{violeta}{azulceleste}%
  {el/la + tio + de + nombre}%
  {mi/tu/su + tio}%
  {es + el tio de + nombre}%
  {tio + tia = tios}%
  {IT: zio | FR: oncle | PT: tio | EN: uncle | CS: stryc | PL: wujek}%
\hspace{1mm}%
\tarjetaPosterior{tia}{violeta}{rojocoral}%
  {el/la + tia + de + nombre}%
  {mi/tu/su + tia}%
  {es + la tia de + nombre}%
  {tio + tia = tios}%
  {IT: zia | FR: tante | PT: tia | EN: aunt | CS: teta | PL: ciotka}%

\vspace{1mm}

\noindent
\tarjetaPosterior{primo}{violeta}{azulceleste}%
  {el/la + primo + de + nombre}%
  {mi/tu/su + primo}%
  {es + el primo de + nombre}%
  {primo + prima = primos}%
  {IT: cugino | FR: cousin | PT: primo | EN: cousin | CS: bratranec | PL: kuzyn}%
\hspace{1mm}%
\tarjetaPosterior{prima}{violeta}{rojocoral}%
  {el/la + prima + de + nombre}%
  {mi/tu/su + prima}%
  {es + la prima de + nombre}%
  {primo + prima = primos}%
  {IT: cugina | FR: cousine | PT: prima | EN: cousin | CS: sestrenice | PL: kuzynka}%

\vspace{1mm}

\noindent
\tarjetaPosterior{marido}{violeta}{azulceleste}%
  {el + marido + de + nombre}%
  {mi/tu/su + marido}%
  {es + el marido de + nombre}%
  {marido y mujer (no marido/marida)}%
  {IT: marito | FR: mari | PT: marido | EN: husband | CS: manzel | PL: maz}%
\hspace{1mm}%
\tarjetaPosterior{mujer}{violeta}{rojocoral}%
  {la + mujer + de + nombre}%
  {mi/tu/su + mujer}%
  {es + la mujer de + nombre}%
  {mujer = woman / wife}%
  {IT: moglie | FR: femme | PT: esposa | EN: wife | CS: manzelka | PL: zona}%

\vspace{1mm}

\noindent
\tarjetaPosterior{nieto}{violeta}{azulceleste}%
  {el/la + nieto + de + nombre}%
  {mi/tu/su + nieto}%
  {es + el nieto de + nombre}%
  {nieto + nieta = nietos}%
  {IT: nipote | FR: petit-fils | PT: neto | EN: grandson | CS: vnuk | PL: wnuk}%
\hspace{1mm}%
\tarjetaPosterior{nieta}{violeta}{rojocoral}%
  {el/la + nieta + de + nombre}%
  {mi/tu/su + nieta}%
  {es + la nieta de + nombre}%
  {nieto + nieta = nietos}%
  {IT: nipote | FR: petite-fille | PT: neta | EN: granddaughter | CS: vnucka | PL: wnuczka}%

\newpage

%% ============================================================
%% CARAS POSTERIORES — Pagina 6 (5 tarjetas restantes)
%% ============================================================

\noindent
\tarjetaPosterior{sobrino}{violeta}{azulceleste}%
  {el/la + sobrino + de + nombre}%
  {mi/tu/su + sobrino}%
  {es + el sobrino de + nombre}%
  {sobrino + sobrina = sobrinos}%
  {IT: nipote | FR: neveu | PT: sobrinho | EN: nephew | CS: synovec | PL: siostrzeniec}%
\hspace{1mm}%
\tarjetaPosterior{sobrina}{violeta}{rojocoral}%
  {el/la + sobrina + de + nombre}%
  {mi/tu/su + sobrina}%
  {es + la sobrina de + nombre}%
  {sobrino + sobrina = sobrinos}%
  {IT: nipote | FR: niece | PT: sobrinha | EN: niece | CS: neter | PL: siostrzenica}%

\vspace{1mm}

\noindent
\tarjetaPosterior{camionero}{naranja}{azulceleste}%
  {es + camionero/camionera}%
  {trabaja + de + camionero}%
  {mi padre/madre + es + camionero}%
  {camionero / camionera}%
  {IT: camionista | FR: camionneur | PT: caminhoneiro | EN: truck driver | CS: ridic | PL: kierowca}%
\hspace{1mm}%
\tarjetaPosterior{enfermera}{naranja}{rojocoral}%
  {es + enfermero/enfermera}%
  {trabaja + de + enfermera}%
  {mi padre/madre + es + enfermera}%
  {enfermero / enfermera}%
  {IT: infermiera | FR: infirmiere | PT: enfermeira | EN: nurse | CS: sestra | PL: pielegniarka}%

\vspace{1mm}

\noindent
\tarjetaPosterior{agricultor}{naranja}{azulceleste}%
  {es + agricultor/agricultora}%
  {trabaja + de + agricultor}%
  {mi padre/madre + es + agricultor}%
  {agricultor / agricultora (no -o/-a)}%
  {IT: agricoltore | FR: agriculteur | PT: agricultor | EN: farmer | CS: zemedelec | PL: rolnik}%

\end{document}
